%!TEX root = ../main.tex
\subsection*{Responses to Reviewer \#2}
\label{cover_letter:r2}

\begin{shaded}
{

\stitle{R2W1:}
\emph{The multi-agent framework uses a limited set of tools for data preparation. It is unclear how the system would perform on datasets requiring more complex or diverse data preparation operations.}

\stitle{R2D3:}
\emph{The operations and operators in AutoPrep seem highly tailored to the datasets used in the experiments. It remains unclear how AutoPrep would perform if the required operations for a new dataset extend beyond its current capabilities. Although a generalization evaluation is included, it is uncertain whether the data issues in the evaluated datasets overlap significantly.}

}
\end{shaded}

\noindent{\bf RESPONSE}: 
We thank the reviewer for highlighting this important issue.
In response, we have added a \textbf{new experiment} to evaluate how \sys performs on a dataset that requires new data prep capabilities (our main change \textbf{M\ref{modify:new-exps} (4)}). 

Specifically, we introduce a new type of data prep operation: \textbf{table transformation}, which was extensively studied in the VLDB 2023 Best Paper~\cite{li2023auto}. The goal of table transformation is to convert a non-relational table format into a relational one. As existing TQA benchmark datasets do not consider this type of data prep issue, we construct a new dataset, \textbf{TransTQ}, with the assistance of LLMs. This dataset is publicly available on our GitHub repository. In particular, we prompt an LLM to identify tables from WikiTQ where inverse transformation operations (\eg pivot and stack~\cite{li2023auto}) can be applied, thereby generating non-relational versions of structured tables. These transformed tables replace the originals to evaluate whether \sys can effectively address such data prep challenges, which differ significantly from the column derivation, normalization, and filtering tasks primarily studied in our work.

%Specifically, since no existing TQA benchmarking dataset presents sufficiently complex data quality issues, which has insignificant overlap with the TQA datasets used in our experiments, we construct a \textbf{new dataset TransTQ} with the assistance of LLMs, which is available on our GitHub repository. Specifically, we prompt an LLM to identify tables from WikiTQ where inverse transformation operations (\eg pivot and stack~\cite{li2023auto}) can be applied. These operations are then executed to generate non-relational table formats, introducing new challenges for TQA on \textbf{table transformation}.

As shown in \textbf{Table~\ref{tbl:new_operator_new_operation} of Section~\ref{subsec:evaluation_on_generalization}}, without explicit guidance for handling table transformations, all baseline TQA methods exhibit poor accuracy. In contrast, thanks to the modular design of our multi-agent framework, incorporating a new logical operation (\term{Transform}) requires only minimal human effort, yet significantly boosts \sys's performance to outperform all baselines, \eg outperforming the second-best method by \textbf{9.31} points. Moreover, we also validate that adding this new operation does not degrade the performance of \sys on the original WikiTQ dataset, confirming that its multi-agent architecture is well extensible to support the seamless integration of new data prep functionalities.

\stitle{Finding: \sys demonstrates strong extensibility to the dataset that requires new logical operation types.}


%Specifically, we have constructed \textbf{a new dataset} TransTQ to evaluate how our \sys performs when ``the required operations for a new dataset extend beyond its current capabilities'' and added \textbf{a new experiment} (Exp~\ref{exp:scalability_exp_new_op_type}) in \textbf{Section~\ref{subsec:evaluation_on_generalization}}. Please refer to \textbf{M\ref{modify:newexp_evaluate_generalizability_on_complex_tqa_dataset}} for further details. The results show that AutoPrep demonstrates \textbf{strong generalizability} to datasets that requires new logical operation types.

% Specifically, considering that there is no appropriate TQA dataset with complicated data quality issues we construct a new Dataset TransTQ from WikiTQ with the assistance of LLMs and publish it in our Github Repo. We then evaluate the generalizability of all methods on the new dataset. Firstly, we select a few (\ie two in our experiments) demos from TransTQ and add the demos to the prompts of all methods. Secondly, we evaluate the performance on WikiTQ after integrating new demos.

% \stitle{How does \sys perform?} Without adding new demos, all methods perform poorly on TransTQ. But after integrating the new demos, \sys achieves new SOTA performance comparing with other methods (improve by \textbf{3.25}), as shown in Table~\ref{tbl:new_operator_new_operation}. Moreover, the performance of \sys does not drop on the original WikiTQ dataset. This indicates that our \sys can generalize well on new data quality issues while maintaining its ability on original issues.

% We further evaluate \sys on datasets beyond its capabilities, \eg a dataset requiring complex transforming operations on non-relational tables in Exp~\ref{exp:scalability_exp_new_op_type}. Considering that there is no TQA dataset with complicated data quality issues, we construct a new Dataset TransTQ from WikiTQ with the assistance of LLMs and publish it in our Github Repo. We compare \sys and other TQA methods on TransTQ with or without integrating ICL demonstrations guiding them to generating transforming operations. We do not compare with CoTable since it depends on a predefined finite-state automaton which determines the search space of the operator and is hard to be modified. \textbf{(2) How \sys performs?} As shown in Table~\ref{tbl:new_operator_new_operation}, without specific guidance of ICL demonstrations, all TQA methods perform poorly, especially \sys. However, when integrating guidance of transforming operations, \sys achieves new SOTA performance comparing with other methods. Recall that we do not claim that \sys is comprehensive in solving all data quality issues, but is scalable with the design of logical-physical operators, which is well demonstrated by the results in Exp~\ref{exp:scalability_exp_new_op_type}.


\begin{shaded}
{
\stitle{R2W2}: 
\emph{Several technical aspects, such as the mechanisms for orchestrating operations and converting analyzer sketches into actions, are not adequately explained in the paper.}

\stitle{R2D1:} 
\emph{The order of operations (and operators) that LLMs orchestrate really matters. However, it does not clarify how this order is determined by the planner and programmer.}

\stitle{R2D2:}
\emph{
I am not sure how the analyzer sketch is converted into operations. In particular, it is unclear how LLMs figure out the purpose of each clause and identify the specific columns to manipulate. I try to find the corresponding prompt in the technical report, but it seems not included.
}
%
}
\end{shaded}

\noindent{\bf RESPONSE}: 
Thanks for the suggestion.
In response, we have revised \textbf{Section~\ref{subsec:planner-cot}} to present more technical details on orchestrating operations and converting analyzer sketches into actions. In addition, we have included all the detailed \textbf{LLM prompts}. See our main change \textbf{M\ref{modify:technical-details}}.

First, for orchestrating operations, given the constructed Analysis Sketch, we orchestrate the logical operations and determine their execution order by translating the sketch into actionable steps.
Specifically, we leverage the syntax parser from Binder~\cite{cheng2023binding} to extract all UDFs and operation clauses from the sketch and sort them according to their intended execution sequence from the parser.
Based on this parsed order, we then generate the corresponding logical operations.

%We have revised \textbf{Section~\ref{subsec:planner-cot}}
% \lei{fix the xxx.}
%to clarify the logical operator generation process.
%When we construct the QA-Sketch, we use the syntax parser from Binder~\cite{cheng2023binding} to extract all UDFs and operation clauses in the sketch and sort them by the execution order. Then, we generate the corresponding logical operators based on the parsed executing order.


% We use the method in Binder~\cite{cheng2023binding} to determine the xxx. We have revised Section xxx to clarify this. 
% \sys parses the generated QA-Sketch to determine the order of logical operator. Afterwards, the \textsc{Programmer} generates the physical operations by the order of logical operators. Specifically, the \textsc{Planner} use a syntax parser to extracts all UDFs and operation clauses in the sketch and sort them by the execution order. We have uploaded a running case illustrating how \sys works in Appendix~\ref{sec:appendix_a_running_example_of_autoprep}.

\input{plots/tex/prompt_instance_generating_logical_requirement}

Second, for converting analyzer sketches into actions, we provide several in-context learning demonstrations that guide the LLM to understand the intent of each clause and identify the corresponding data preparation requirements.
%
For better illustration, we provide an example in Figure~\ref{fig:prompt_instance_generating_logical_requirement}. In the prompt, we begin by defining the task through a description and an illustrative example, which includes both the input and the desired output. This helps the LLM understand that the objective is to generate the requirement for a column derivation operation. Next, we present the source columns for analysis and specify that the goal is to derive the column \term{is\_loss} in the Sketch to answer question Q. With this information, the LLM is guided to infer the intent of each clause and identify the corresponding data preparation requirements. Additional details of the prompt construction are provided in our technical report~\cite{technical_report}.

%For better illustration, we provide an example as shown in Figure~\ref{fig:prompt_instance_generating_logical_requirement}. In the prompt, we first define the task by providing the task description and an illustrative example consisting of the input and desired output, from which the LLM can learn the task is to produce the requirement of a column derivation operation. Then, we provide the source columns for the LLM to analyze and remind it that the requirement is to derive the column `is\_loss' in the QA-Sketch to answer question Q. Finally, with all the information provided, the LLM can figure out the purpose of each clause and identify the specific data prep requirements to satisfy.
%We have presented more details of the prompts in our technical report~\cite{technical_report}.

%We have revised \textbf{Section~\ref{subsec:planner-cot}} and added detailed prompts
% (Figure~\ref{fig:prompt_planner_generate_req_aug} and Figure~\ref{fig:prompt_planner_generate_req_norm})
%in our \textbf{Technical Report}.

\begin{shaded}
{
\stitle{R2W3:}
\emph{AutoTQA studies a similar problem, which reportedly achieves better results. A detailed comparison and discussion of the differences between AutoPrep and AutoTQA are necessary.}

\stitle{R2D4:}
\emph{AutoTQA studies a similar problem, which reportedly achieves better performance. The paper does not include a comparative study with AutoTQA, nor does it discuss the key differences between AutoTQA and AutoPrep.}

}
\end{shaded}

\noindent{\bf RESPONSE}: 
Thanks for the suggestion. 
We have added a \textbf{new experiment} to compare with AutoTQA (our main change \textbf{M\ref{modify:new-exps} (2)}), and discussed the differences between AutoPrep and AutoTQA.
The experimental results show that \sys outperforms AutoTQA, which also adopts a multi-agent, LLM-based framework.
Please refer to the response to \textbf{R1D13} for detailed results and analysis. 

%We have added this baseline in the revised version. Please refer to \textbf{M\ref{modify:more_baselines}} for more details and refer to \textbf{R1D13} for the comparison between \sys and AutoTQA.

\begin{shaded}
{
\stitle{R2D5:}
\emph{Discuss the limitations of AutoPrep.}
}
\end{shaded}

\noindent{\bf RESPONSE}: 
We have discussed the limitations of \sys, which provide promising research directions, in the \textbf{Section~\ref{sec:conclusion_and_future_work} Conclusion and Future Work} (our main change \textbf{M\ref{modify:future-work}}). 

First, we aim to extend the system's current data prep capabilities by integrating a broader set of operations to handle more complex issues such as missing values, duplicates, etc. Second, \sys currently focuses on single-table question answering; extending its capabilities to support multi-table TQA is essential for tackling more realistic and complex scenarios.
The third direction is to enable question-aware data preparation over enterprise datasets, which are much more complex than existing TQA benchmarks. 
% enabling question-aware data preparation over enterprise datasets.

%\textbf{Section~\ref{sec:conclusion_and_future_work}}. 

%\stitle{\sys only supports QA on the single table.} In practical scenarios, it is possible that answering a question correctly needs to access the data from different tables. More complicated data quality issues could occur when integrating the tables. For example, how to solve the conflicting fact in two different tables? In addition, data inconsistency issues could occur more frequently, since different tables may come from different sources.
%
%\stitle{\sys only supports QA on the known table.} In open-domain question answering, users will only provide a question and require the system to retrieve relevant data (\eg tables) to answer it. Supporting open QA is one of our future research directions.

%We will extend the capabilities of \sys according to the two limitations above in the future work.










% \begin{shaded}
% {
% \noindent \emph{\underline{R2D1}: 
% % Further explain the order of operation generation.
% D1: The order of operations (and operators) that LLMs orchestrate really matters. However, it does not clarify how this order is determined by the planner and programmer.
% }
% }
% \end{shaded}
% \vspace{-0.5em}
% \noindent{\bf RESPONSE}: \textbf{M\ref{modify:new-exps}-\ref{M\ref{modify:new-exps}:describe_coc}}



% \begin{shaded}
% {
% \noindent \emph{\underline{R2D2}: 
% % Further describe the logical operation generation process
% D2: I am not sure how the analyzer sketch is converted into operations. In particular, it is unclear how LLMs figure out the purpose of each clause and identify the specific columns to manipulate. I try to find the corresponding prompt in the technical report, but it seems not included.
% }
% }
% \end{shaded}
% \vspace{-0.5em}
% \noindent{\bf RESPONSE}: \textbf{M\ref{modify:new-exps}-\ref{M\ref{modify:new-exps}:describe_coc}} 

% We scan the generated QA sketch to extract all operations on columns and supervise LLM with ICL demonstrations to generate logical op. We add the prompt in the technical report.

% % 通过ICL给例子，让大模型直接生成
% % 我们依次扫描qa sketch，某个列上的操作，我们询问LLM给出logical op
% % 添加prompt到technical report


% \begin{shaded}
% {
% \noindent \emph{\underline{R2D3}: 
% % Further generalization evaluation.
% D3: (1) The operations and operators in AutoPrep seem highly tailored to the datasets used in the experiments. It remains unclear how AutoPrep would perform if the required operations for a new dataset extend beyond its current capabilities. (2) Although a generalization evaluation is included, it is uncertain whether the data issues in the evaluated datasets overlap significantly.
% }
% }
% \end{shaded}
% \vspace{-0.5em}
% \noindent{\bf RESPONSE}: We construct a Transform dataset with pivot and stack operators from AutoTable~\cite{li2023auto}. We first utilizing a powerful LLM Deepseek-v3 to select all possible tables to be pivoted or stacked and provide corresponding transforming suggestions. Then, we manually transform the table to a structure not suitable for answering the question. The Transform dataset contains 241 instances in total with 41 stacking instances and 200 pivoting instances. The dataset can be accessed in our Github Repo. We evaluate \sys and other TQA methods which are easy to be integrating to new operations. Notice that we do not compare CoTable~\cite{wang2024chain} because it define a finite state automata to reduce the search space of existing ops. Thus, a new operator Transform can not be plugged in.

% As shown in Table~\ref{tbl:new_operator_new_operation}, after integrating Transform operations, \sys achieves SOTA performance. Moreover, the performance on origin dataset WikiTQ has been improved slightly. \fmh{explain after i analyze the exp result in depth.}

% % First, we define what ``is a new dataset extend beyond its current capabilities''. Since the abilities of \sys is decided by LLM's knowledge and defined operators, the out-of-scope data quality can be categorized as:

% % \begin{itemize}
% %     \item \textbf{Within LLM's knowledge, out of defined operators.}
% %     \item \textbf{Out of LLM's knowledge, within defined operators.} For example, some question will requires domain-specific calculation to generate new columns. The challenges of this problem is to inject the domain knowledge into LLM, which can be solved by continual training or RAG. This is not the focus of our paper.
% %     \item \textbf{Out of LLM's knowledge, Out of defined operators.} To solve this, we need to solve the above two types of questions.
% % \end{itemize}

% % Thus, we need to evaluate \sys on problems within LLM's knowledge, out of defined operators. We design two scenario

% % \sys can be well adapted to this types of problems because: (1) we define two types of operations to solve the problems of incomplete physical operations, i.e., use LLM-generated code for questions failing on defined operations and use LLM prediction for questions can not be solved by code. Refer to R1D4. (2) The Logical-physical design is easy to be extended. We add description of specific extending procedures in the paper. In addition, we design experiment of new operator and new operation, showing that with simple extending effort, \sys improves obviously.

% Seen in Table~\ref{tbl:new_operator_new_operation}

% ``the required operations for a new dataset extend beyond its current capabilities.''

% \input{tables/new_operator_new_operation}

% \begin{shaded}
% {
% \noindent \emph{\underline{R2D4}: 
% % Comparison with AutoTQA
% D4: AutoTQA studies a similar problem, which reportedly achieves better performance. The paper does not include a comparative study with AutoTQA, nor does it discuss the key differences between AutoTQA and AutoPrep.
% }
% }
% \end{shaded}
% \vspace{-0.5em}
% \noindent{\bf RESPONSE}: Compare with AutoTQA.